\chapter{Bloating}

\section*{Definition}
Abdominal bloating is a feeling of fullness, pressure, or enlargement; abdominal distension is visible or measurable increase in abdominal size. Bloating does not always mean that excess gas is present.

\section*{When to See a Doctor}

\begin{itemize}
 \item Arrange assessment for persistent or recurrent bloating, particularly with pain, vomiting, diarrhea or constipation, unexplained weight loss, anemia, bleeding, or a new change in bowel habits.
 \item Seek urgent care for severe or rapidly worsening pain, repeated vomiting, a swollen rigid abdomen, inability to pass stool or gas, blood in the stool, faintness, or fever.
\end{itemize}

\section*{How It Develops}
Gas accumulation may result from swallowed air, bacterial fermentation, or impaired gas transit. Visceral hypersensitivity can amplify the perception of normal gas volumes in people with functional bowel disorders.

\section*{Causes}
\begin{itemize}
 \item Functional gastrointestinal disorders such as irritable bowel syndrome
 \item Small intestinal bacterial overgrowth in selected people with risk factors
 \item Constipation
 \item Carbohydrate intolerance or sensitivity, including lactose or fructose; do not assume that gluten is the cause without appropriate celiac testing
 \item Swallowing air (aerophagia)
 \item Gastroparesis or delayed gastric emptying
 \item Ovarian pathology (cysts, masses)
 \item Celiac disease
 \item Gastrointestinal obstruction
 \item Ascites from liver disease or malignancy
\end{itemize}

\section*{Related Symptoms}
\begin{itemize}
 \item Abdominal pain
 \item Flatulence
 \item Constipation
 \item Diarrhea
 \item Nausea
\end{itemize}

\section*{Medical Evaluation}
\begin{itemize}
 \item Your doctor will take a dietary history and perform an abdominal examination assessing for distension, tympany, fluid wave (ascites), and palpable masses.
 \item Breath testing for lactose intolerance, fructose malabsorption, or small intestinal bacterial overgrowth (SIBO) may be arranged based on symptom pattern.
 \item Abdominal imaging or endoscopy is reserved for alarm features, recent worsening, or an abnormal examination; the test depends on the suspected cause.
 \item Celiac serology may be appropriate, and testing should generally occur while the person is still eating gluten. Breath tests are used selectively for carbohydrate intolerance or SIBO.
\end{itemize}

\section*{Treatment Options}
\begin{itemize}
 \item Dietary changes should be individualized. A time-limited low-FODMAP approach may help some people with a functional disorder, preferably with a dietitian and followed by systematic reintroduction; restrictive diets can cause nutritional problems.
 \item Do not start a gluten-free diet before celiac testing unless a clinician advises it. Antibiotics such as rifaximin are reserved for selected diagnosed conditions and are not routine treatment for bloating or a positive breath test alone.
 \item Prokinetic agents (metoclopramide, domperidone) or low-dose tricyclic antidepressants may be prescribed for gastroparesis or visceral hypersensitivity.
 \item Treatment targets the identified cause, such as an appropriate bowel regimen for constipation or specialist management of gastroparesis, pancreatic disease, liver disease, or ascites.
\end{itemize}

\section*{Self-Care and Home Management}
\begin{itemize}
 \item Eat smaller, more frequent meals rather than large portions to reduce gastric distension and gas production.
 \item Keep a food-symptom diary to identify trigger foods (common culprits include beans, lentils, cabbage, onions, carbonated drinks, and artificial sweeteners).
 \item A short trial of an over-the-counter product may help some people, but discuss persistent symptoms with a clinician. Use a low-FODMAP diet only with guidance and a plan to reintroduce foods.
 \item Eat slowly, chew food thoroughly, and avoid drinking through straws or chewing gum to reduce aerophagia (swallowed air).
\end{itemize}

\section*{References}
\begin{thebibliography}{99}
\bibitem{bloating:ref1} Lacy BE, Mearin F, et al. ``Bowel disorders.'' \textit{Gastroenterology}. 2016;150(6):1393-1407.
\bibitem{bloating:ref2} Azpiroz F, Malagelada JR. ``Abdominal bloating.'' \textit{Gastroenterology}. 2005;129(3):1060-1078.
\bibitem{bloating:ref3} Ford AC, Moayyedi P, et al. ``American College of Gastroenterology monograph on the management of IBS and chronic idiopathic constipation.'' \textit{Am J Gastroenterol}. 2014;109(Suppl 1):S2-S26.
\bibitem{bloating:ref4} Feldman M, Friedman LS, Brandt LJ, eds. \textit{Sleisenger and Fordtran\'s Gastrointestinal and Liver Disease}, 11th ed. Elsevier, 2021.
\bibitem{bloating:ref5} American Gastroenterological Association. ``Evaluation and management of belching, abdominal bloating and distention.'' Clinical Practice Update, 2023. https://gastro.org/clinical-guidance/evaluation-and-management-of-belching-abdominal-bloating-and-distention/.
\end{thebibliography}
